\documentclass[french,a4paper]{article}

\usepackage[utf8]{inputenc}
\usepackage[T1]{fontenc}
\usepackage{lmodern}
\usepackage{geometry}
\usepackage{graphicx}
\usepackage{babel}
\usepackage{amsmath}
\usepackage{amsthm}
\usepackage{amssymb}
\usepackage{hyperref}
\usepackage{stmaryrd}
\usepackage{enumitem}
\usepackage{tcolorbox}
\usepackage{dsfont}
\usepackage{multirow}
\usepackage{etoc}
\usepackage{titlesec}
\usepackage{esint}
\usepackage{cancel}
\usepackage{mathdots}

\setlist[itemize]{label=$\bullet$}


\renewcommand{\P}{\mathbb{P}}
\newcommand{\N}{\mathbb{N}}
\newcommand{\Z}{\mathbb{Z}}
\newcommand{\Q}{\mathbb{Q}}
\newcommand{\R}{\mathbb{R}}
\newcommand{\C}{\mathbb{C}}
\newcommand{\K}{\mathbb{K}}
\renewcommand{\L}{\mathbb{L}}
\newcommand{\U}{\mathbb{U}}
\newcommand{\st}{^*}
\newcommand{\tq}{\middle|}
\newcommand{\inv}{^{-1}}
\newcommand{\E}{\mathbb{E}}
\newcommand{\F}{\mathbb{F}}
\newcommand{\V}{\mathbb{V}}
\renewcommand{\ker}{\text{Ker}}
\newcommand{\im}{\text{Im}}
\newcommand{\card}{\text{Card}}
\newcommand{\Tr}{\text{Tr}}

\theoremstyle{plain}
\newtheorem{thm}{Théorème}
\newtheorem*{ind}{Indication}

\theoremstyle{definition}
\newtheorem{dfn}{Definition}

\theoremstyle{remark}
\newtheorem*{rmq}{Remarque}
\newtheorem*{app}{Application}

\newtcolorbox[auto counter]{ex}[2][]{colback=white, colframe=green!50!white, coltitle=black, fonttitle=\bfseries, title={Exercice~\thetcbcounter~: #2}, after title={\hfill #1}}

\title{TD Intégrales généralisées}

\begin{document}

\maketitle

\section{Intégration sur un segment}
\begin{ex}{$\Delta\star$ Formule de la moyenne}
	Soit $f$ et $g$ deux applications continues de $[a,b]$ avec $a<b$. On suppose que $\forall t\in[a,b],\ g(t)\geqslant0$. Alors : $$\exists c\in[a,b],\ \int_{a}^{b}f(t)g(t)\ dt=f(c)\int_{a}^{b}g(t)\ dt$$
\end{ex}
\begin{proof}
	$f$ est continue sur le compact $[a,b]$ donc d'après le théorème des bornes atteintes, il existe $(\alpha,\beta)\in[a,b]^2$ tel que : $$\forall t\in[a,b],\ f(\alpha)\leqslant f(t)\leqslant f(\beta)$$ Par positivité de $g$, on a : $$\forall t\in[a,b],\ f(\alpha)g(t)\leqslant f(t)g(t)\leqslant f(\beta)g(t)$$ Par croissance de l'intégrale et linéarité de l'intégrale, on en déduit que $$f(\alpha)\int_{a}^{b}g(t)\ dt \leqslant \int_{a}^{b}f(t)g(t)\ dt\leqslant f(\beta)\int_{a}^{b}g(t)\ dt$$ Si $\displaystyle \int_{a}^{b}g(t)\ dt=0$, alors il suffit de prendre $c=a$ car tous les termes de l'inégalité précédente sont nuls. Sinon, $\displaystyle \int_{a}^{b}g(t)\ dt >0$, si bien que l'encadrement obtenu précédemment se réécrit : $$f(\alpha)\leqslant\displaystyle\frac{\int_{a}^{b}f(t)g(t)\ dt}{\int_{a}^{b}g(t)\ dt}\leqslant f(\beta)$$ Le théorème des valeurs intermédiaires permet alors de conclure.
\end{proof}
\begin{rmq}
	La preuve s'adapte dans le cas où $g$ est négative.
\end{rmq}
\begin{app}
	Permet de déduire l'égalité de Taylor-Lagrange de la formule de Taylor avec reste intégrale lorsque $f$ est de classe $\mathcal{C}^{n+1}$.
\end{app}
\begin{ex}{Où sont les intégrales ?}
	Montrer que $2<\pi<4$.
\end{ex}
\begin{proof}
	On note $f:x \mapsto \frac{1}{1+x^2}$. On remarque que sur sur le segment $[0,1]$, on a $\frac{1}{2} < f <1$. Il suffit alors d'intégrer sur ce segment avec la stricte croissance de l'intégrale, puis de multiplier le tout par 4. On peut aussi raisonner avec un quart de cercle ou un demi-cercle dans la plan, mais c'est alors un peu plus embêtant de justifier les inégalités.
\end{proof}

\begin{ex}{$\Delta$ Avec la fonction réciproque}
	Soit $f:[0,1]\to[0,1]$ une bijection strictement croissante de classe $\mathcal{C}^1$. Montrer que $$\int_{0}^{1}f+\int_{0}^{1}f\inv=1$$ Interpréter géométriquement.
\end{ex}

\begin{proof}
	Faire un changement de variable dans la deuxième. Le résultat se comprend bien car la courbe de la bijection réciproque est le symétrique de la courbe par rapport à la première bissectrice du repère.
\end{proof}

\begin{ex}{$\Delta$ Théorème des moments}
	Soit $f\in\mathcal{C}^0([a,b],\R)$ et $n\in\N$. Montrer que si $$\forall i\in\llbracket0,n\rrbracket,\ \int_{a}^{b}t^if(t)\ dt=0$$ alors $f$ s'annule au moins $n+1$ fois sur $]a,b[$.
\end{ex}

\begin{proof}
	A FAIRE
\end{proof}

\begin{ex}{$\star\star$ Difficile}
	Déterminer les polynômes $P$ qui envoient $[0,1]$ dans lui-même et tels que $$\forall f\in\mathcal{C}^0([0,1],\R),\ \int_{0}^{1}f(P(t))\ dt=\int_{0}^{1}f(t)\ dt$$
\end{ex}

\begin{proof}
	Ce ne sont que $X$ et $1-X$. Il faut (je crois) prendre des fonctions qui s'approchent de Diracs pour le prouver.
\end{proof}
\section{Suites d'intégrales sur un segment}

\begin{ex}{Une limite}
	Déterminer $$\lim_{{n\to+\infty}}\int_{0}^{\pi/2}(\sin(x))^{1/n}\ dx$$
\end{ex}

\begin{proof}
	Utiliser l'inégalité :$$\forall x\in\left[0,\frac{\pi}{2}\right],\ \frac{2}{\pi}x\leqslant\sin(x)\leqslant x$$ pour se ramener à des intégrales de fonctions puissances. Calculer les intégrales qui encadrent et déterminer leur limite. Le théorème des gendarmes permet de conclure que la limite recherchée est $\displaystyle\frac{\pi}{2}$.
\end{proof}

\begin{ex}{$\Delta$ Norme infinie}
	Soit $f\in\mathcal{C}^0([a,b],\R_+)$ avec $a<b$. Déterminer $$\lim_{n\to+\infty}\left(\int_{a}^{b}f(t)^{n}\ dt\right)^{1/n}$$
\end{ex}

\begin{proof}
	Il s'agit de $\lVert f\rVert_\infty$. Il faut le faire par encadrement avec des $\varepsilon$. Un côté de l'inégalité est simple. Pour l'autre, s'approcher à $\eta$ de là où on est à $\varepsilon$ du maximum puis minorer brutalement le reste. On s'en sort.
\end{proof}

\begin{ex}{$\star$ Lemme de Riemann-Lebesgue sur un segment}
	Soit $f$ une fonction continue par morceaux sur un segment $[a,b]$ avec $a<b$.
	\begin{enumerate}
		\item Déterminer $$\lim_{\lvert\lambda\rvert\to+\infty}\int_{a}^{b}f(t)e^{i\lambda t}\ dt$$ On pourra commencer par le cas où $f$ est de classe $\mathcal{C}^1$, puis en escalier.
		\item Que devient le résultat si on remplace $e^{i\lambda t}$ par $g(\lambda t)$ où $g$ est une fonction continue périodique. On pourra commencer par le cas où $g$ est de valeur moyenne nulle.
	\end{enumerate}
\end{ex}

\begin{proof}
	\begin{enumerate}
		\item Dans le cas $\mathcal{C}^1$, une IPP en primitivant l'exponentielle suffit. Le cas d'une constante est un cas particulier du cas $\mathcal{C}^1$. Dans le cas en escalier, on prend une subdivision adaptée pour se ramener à un nombre fini de fonctions en escalier. Dans le cas général, on peut approximer $f$ par une fonction en escalier puis on travaille avec du $\varepsilon$.
		\item Si $g$ est continue périodique de valeur moyenne nulle, ses primitives sont aussi périodiques donc bornées et le résultat est le même que précédemment. Dans le cas général, on peut écrire $g=\langle g\rangle+g_0$ où $g_0$ est donc continue périodique de valeur moyenne nulle. La limite recherchée est donc $$\langle g\rangle\int_{a}^{b}f(t)\ dt$$
	\end{enumerate}
\end{proof}

\begin{ex}{$\star$ Irrationalité de $\pi$}
	On se propose de montrer par l'absurde que $\pi$ est irrationnel. Pour cela, on suppose $\pi=a/b$ avec $a$ et $b$ dans $\N\st$ et l'on pose $$I_n=\int_{0}^{\pi}\frac{x^n(a-bx)^n}{n!}\sin(x)\ dx$$
	\begin{enumerate}
		\item Montrer que $(I_n)$ converge vers $0$.
		\item Montrer que toutes les dérivées de $\displaystyle \frac{X^n(a-bX)^n}{n!}$ sont entière en $0$ puis en $\pi$
		\item En déduire que $I_n$ est un entier puis conclure.
	\end{enumerate}
\end{ex}

\begin{proof}
	Grand classique ! Pour les détails, voir la pâle de sup.
	\begin{enumerate}
		\item Majorer le numérateur par une constante à la puissance $n$. La factorielle l'emporte.
		\item Remarquer qu'au-delà d'un certain degré, toutes les dérivées sont nulles. Pour les autres : remarquer que $0$ est racine de multiplicité $n$ donc les premières dérivées sont nulles. Pour $\pi$, remarquer que le polynôme est symétrique ($P(\pi-X)=P(X)$).
		\item Faire des intégrations par parties successives (mieux vaut utiliser la formule de la partie astuces pour ne pas se perdre !). Ensuite, une suite d'entier qui tend vers $0$ est stationnaire. Mais les fonctions qu'on intègre sont continues positives non identiquement nulles, donc cela est absurde.
	\end{enumerate}
\end{proof}

\section{Calcul intégral sur un intervalle quelconque}
\begin{ex}{Convergence et calcul}
	Convergence et calcul éventuel des intégrales : $$\int_{0}^{1}\frac{x}{\sqrt{1-x}}\ dx \ \ \text{et} \ \ \int_{-1}^{1}\frac{dx}{\sqrt{x^2(1+x)}}$$
\end{ex}
\begin{proof}
	On commence par vérifier la continuité par morceaux sur ! Pour la première, une primitive fournit l'existence et le calcul (=1). Et on n'a pas de continuité par morceaux en $0$ pour la deuxième et pas de prolongement possible (non bornée au voisinage de $0$) donc l'intégrale n'existe pas.
\end{proof}
\begin{ex}{Avec du cosinus hyperbolique}
	Étant donné un réel $\alpha$, convergence et calcul de $$I(\alpha)=\int_{0}^{+\infty}\frac{dt}{\cosh(t)+\cosh(\alpha)}$$
\end{ex}
\begin{proof}
	On a continuité (par morceaux) sur $[0,+\infty[$ et un équivalent en $\displaystyle\frac{2}{e^t}$ en $+\infty$ donc l'intégrabilité et la convergence de l'intégrale quel que soit $\alpha$. Ensuite, remplacer $\cosh(t)$ par des exponentielles et multiplier le dénominateur et le numérateur par $e^t$. Faire un changement de variable $u=e^t$. On obtient : $$I(\alpha)=\int_{1}^{+\infty}\frac{2 \ du}{u^2+2u\cosh(\alpha)+1}$$ Ensuite, distinguer si $\alpha =0$ (on trouve $1$) ou non : dans ce dernier cas, on trouve après factorisation et DES $$I(\alpha)=\frac{1}{\sinh(\alpha)}\ln\left(\frac{1+e^\alpha}{1+e^{-\alpha}}\right)$$
\end{proof}
\begin{ex}{Partie fractionnaire et inverse}
	Montrer que la fonction $f$ définie sur $]0,1]$ par : $$f(t)=\frac{1}{t}-\left\lfloor\frac{1}{t}\right\rfloor$$ est intégrable et calculer son intégrale.
\end{ex}
\begin{proof}
	On a une continuité par morceaux due au caractère de partie fractionnaire. Ensuite, poser $$u_k=\int_{1/(k+1)}^{1/k}f(t)\ dt$$ et considérer les sommes partielles. Calculer sur ces intervalles-là pour faire apparaître dans la somme partielle une série harmonique, puis utiliser le développement asymptotique habituel avec la constante d'Euler $\gamma$. Cela fournit la convergence et la valeur : $$\int_{0}^{1}f(t)\ dt=1-\gamma$$
\end{proof}
\begin{ex}{$\Delta\star$ Intégrale de Dirichlet}
	Convergence et calcul de l'intégrale de Dirichlet $$I=\int_{0}^{\pi/2}\ln(\sin(x))\ dx$$
\end{ex}
\begin{proof}
	La convergence se fait bien. Ne pas oublier la continuité-par-morceaux-sur ! Montrer que cela ne change rien de transformer le $\sin$ en $\cos$ avec le changement de variable $u=\displaystyle\frac{\pi}{2}-x$. Ensuite, calculer $2I$ et utiliser la propriété fonctionnelle du $\ln$ : $$2I=-\frac{\pi\ln(2)}{2}+\underbrace{\int_{0}^{\pi/2}\ln(\sin(2x))\ dx}_{=J}$$ Or, un changement de variable $u=\displaystyle\frac{x}{2}$ puis un découpage en 2 et l'utilisation du fait que $\sin(\pi-u)=\sin(u)$ donnent $J=I$. Miraculeusement, cela se termine et on obtient alors : $$I=-\frac{\pi\ln(2)}{2}$$ 
\end{proof}
\begin{ex}{$\star$ Beurk}
	Convergence et calcul de $\displaystyle\int_{0}^{+\infty}\frac{\ln(1+t^2)}{1+t^2}\ dt$ et $\displaystyle\int_{0}^{+\infty}\left(\frac{\arctan(t)}{t}\right)^2 \ dt$.
\end{ex}
\begin{proof}
	Pour la première, on effectue le changement de variable $t=\tan(\theta)$ pour se ramener à -2 fois l'intégrale de Dirichlet. On obtient alors $\pi\ln(2)$. Pour la deuxième, même changement de variable puis 2 IPP pour se ramener encore une fois à -2 fois l'intégrale de Dirichlet.
\end{proof}
\begin{rmq}
	Askip il est possible de le faire en utilisant l'exercice suivant.
\end{rmq}
\begin{ex}{Utilisation du changement de variable inverse}
	Montrer que $\displaystyle I=\int_{0}^{+\infty}\frac{\ln(t)}{1+t^2}\ dt=0$ et calculer $$J_a=\int_{1/a}^{a}\frac{\arctan(t)}{t}\ dt$$ pour $a>0$.
\end{ex}
\begin{proof}
	Le changement de variable $u=1/t$ donne $I=-I$ donc $I=0$. Pour $J_a$, faire une IPP et on tombe sur une deuxième intégrale $$\int_{1/a}^{a}\frac{\ln(t)}{1+t^2}\ dt$$ Le même changement de variable montre que celle-ci est encore nulle. Avec le terme tout intégré, on a donc : $$J_a=\frac{\pi}{2}\ln(a)$$
\end{proof}
\section{Intégration sur un intervalle quelconque}
\begin{ex}{$\star$ Intégrale de l'écart à la translation}
	Soit $f:R\to\C$ intégrable. Déterminer : $$\lim_{x\to+\infty}\int_{-\infty}^{+\infty}\lvert f(t-x)-f(t)\rvert\ dt$$ Indication : on pourra se ramener à $\displaystyle\int_{-\infty}^{+\infty}\lvert f(t-x)-f(t+x)\rvert\ dt$
\end{ex}
\begin{proof}
	Pour se ramener à l'indication, faire un changement de variable pour moyenner. Il faudra que je fasse la suite un jour...
\end{proof}
\begin{ex}{Intégrabilité et uniforme continuité}
	Soit $f$ uniformément continue sur $\R_+$ et intégrable sur $\R_+$.
	\begin{enumerate}
		\item On suppose $f$ réelle positive. Montrer que $f$ converge vers $0$ en $+\infty$.
		\item Que peut-on dire si $f$ est à valeurs réelles quelconques ? à valeurs complexes ?
		\item Que deviennent ces résultats si l'on remplace l'hypothèse d'intégrabilité par la convergence de l'intégrale $\displaystyle\int_{0}^{+\infty}f(t)\ dt$ ?
		\item Et si l'on suppose seulement $f$ continue ? 
	\end{enumerate}
\end{ex}
\begin{proof}
	A faire un jour...
\end{proof}
\begin{ex}{$\Delta$ Intégrable positive décroissante}
	Soit $f:\R_+\to\R_+$ continue par morceaux et décroissante.
	\begin{enumerate}
		\item Montrer que si $f$ est intégrable sur $\R_+$, alors $\lim_{t\to+\infty}tf(t)=0$. Étudier la réciproque.
		\item Le résultat subsiste-t-il si l'on ne suppose plus $f$ décroissante ? 
	\end{enumerate}
\end{ex}
\begin{proof}
	\begin{enumerate}
		\item $f$ admet une limite (théorème de la limite monotone) qui est nulle sans quoi $f$ n'est pas intégrable par intégration des relations de comparaison. Ensuite, en majorant par le reste, on a : $$\int_{x/2}^{x}f(t)\ dt\to 0$$ La décroissance fournit donc $\displaystyle\frac{x}{2}f(x)\to 0$ et le résultat recherché. La réciproque est fausse, il suffit de considérer une intégrale de Bertrand avec la fonction $$x\mapsto\frac{1}{x\sqrt{\ln(x)}}$$
		\item Le résultat est alors faux. On peut prendre $f$ nulle presque partout et lui faire faire ce que l'on souhaite ailleurs (avec des pics de plus en plus fins).
	\end{enumerate}
\end{proof}
\begin{ex}{$\Delta$ Inégalité de Hardy}
	Soit $f:\R_+\to\R$ continue de carré intégrable. On pose $$F(x)=\frac{1}{x}\int_{0}^{x}f(t)\ dt$$
	\begin{enumerate}
		\item Montrer que $F(x)$ est négligeable devant $1/\sqrt{x}$ au voisinage de $+\infty$.
		\item $\star$ Montrer que $F$ est de carré intégrable sur $\R_+\st$. Et sur $\R_+$ ?
		\item $\star$ Trouver la plus petite constante $C>0$ possible qui vérifie l'\textbf{inégalité de Hardy} : $$\int_{0}^{+\infty}F(x)^2\ dx\leqslant C\int_{0}^{+\infty}f(x)^2\ dx$$
	\end{enumerate}
\end{ex}
\begin{proof}
	\begin{enumerate}
		\item L'inégalité de Cauchy-Schwarz fournit un grand O. Retravailler ce grand O en découpant pour obtenir le petit o.
		\item Remarquer que $F$ est dérivable que que pour tout $x>0$, on a $$F'(x)=\frac{f(x)-F(x)}{x}$$ Intégrer par parties $F^2$ entre $0$ et $x$ en dérivant $F^2$. On obtient $$\int_{0}^{x}F^2(t)\ dt=2\int_{0}^{x}f(t)F(t)\ dt-F^2(x)$$ Dégager le dernier terme qui est négatif, puis utiliser l'inégalité de Cauchy-Schwarz en distinguant le cas où $F^2$ est identiquement nulle. On obtient alors $$\int_{0}^{x}F^2(t)\ dt\leqslant4\int_{0}^{x}f^2(t)\ dt\leqslant4\int_{0}^{+\infty}f^2(t)\ dt$$ On obtient bien l'intégrabilité souhaitée.
		\item On passe à la limite dans l'inégalité précédente. La constante $4$ est optimale, le vérifier avec une fonction simple comme l'exponentielle.
	\end{enumerate}
\end{proof}
\begin{ex}{$\star\star$ Utilisation d'une formule de Taylor}
	Soit $f\in\mathcal{C}^2(\R_+,\R)$ intégrable telle que $f''$ soit de carré intégrable. Montrer que $f$ et $f'$ admettent une limite nulle en $+\infty$. On pourra utiliser une formule de Taylor.
\end{ex}
\begin{proof}
	On prouve rapidement en passant à la valeur absolue et en majorant par un reste que $$\int_{x}^{x+1}f(t)\ dt\xrightarrow[x\to+\infty]{}0$$ On prend $F$ une primitive de $f$. On applique la formule de Taylor avec reste intégrale à $F$ qui est de classe $\mathcal{C}^3$ : $$F(x+1)-\frac{1}{0!}F(x)=\frac{1}{1!}f(x)+\frac{1}{2}f'(x)+\int_{x}^{x+1}\frac{(x+1-t)^2}{2!}f''(t)\ dt$$ Or, l'inégalité de Cauchy-Schwarz montre que la dernière intégrale tend vers $0$ en utilisant l'hypothèse. Donc $$f(x)+\frac{1}{2}f'(x)\xrightarrow[x\to+\infty]{}0$$ Et là il faut conclure ... (soit j'ai mal noté, soit Pierre n'avait pas réellement fini l'exercice).
\end{proof}
\section{Intégrales semi-convergentes}
\begin{ex}{$\Delta\star$ Intégrale de Frullani}
	Soit $f$ continue de $\R_+$ dans $\C$ admettant une limite finie $l$ en $+\infty$.
	\begin{enumerate}
		\item Pour $b>a>0$, montrer la convergence de $I_{a,b}=\displaystyle\int_{0}^{+\infty}\frac{f(bx)-f(ax)}{x}\ dx$ et la calculer.
		\item Application : calcul de $\displaystyle\int_{0}^{1}\frac{x-1}{\ln(x)}\ dx$
	\end{enumerate}
\end{ex}
\begin{proof}
	\begin{enumerate}
		\item On prend $0<u<v$. On calcule : \begin{align*}
			\int_{u}^{v}\frac{f(bx)-f(ax)}{x}\ dx&=\int_{u}^{v}\frac{f(bx)}{x}\ dx-\int_{u}^{v}\frac{f(ax)}{x}\ dx \\
			&= \int_{bu}^{bv}\frac{f(t)}{t}\ dt-\int_{au}^{av}\frac{f(t)}{t}\ dt \\
			&=\int_{av}^{bv}\frac{f(t)}{t}\ dt-\int_{au}^{bu}\frac{f(t)}{t}\ dt \\
		\end{align*} avec des changements de variables et en découpant avec la 	relation de Chasles. Ensuite : $$\int_{av}^{bv}\frac{f(t)}{t}\ dt=\underbrace{\int_{av}^{bv}\frac{f(t)-l}{t}\ dt}_{\to0}+l\ln\left(\frac{b}{a}\right)$$ De même, on trouve : $$\int_{au}^{bu}\frac{f(t)}{t}\ dt\xrightarrow[u\to 0]{}f(0)\ln\left(\frac{b}{a}\right)$$ Ainsi, l'intégrale $I_{a,b}$ converge et on a : $$I_{a,b}=(l-f(0))\ln\left(\frac{b}{a}\right)$$
		\item Effectuer le changement de variable $x=e^{t}$ et utiliser la question précédente pour obtenir $\ln(2)$.
	\end{enumerate}
\end{proof}
\begin{ex}{Une semi-convergence}
	Convergence de $\displaystyle\int_{0}^{+\infty}\frac{\sin(t)}{\sqrt{t+\cos(t)}}\ dt$
\end{ex}
\begin{proof}
	Pas de problème de continuité par morceaux, de définition ni au voisinage de $0$. Faire un développement asymptotique au voisinage de $+\infty$ pour obtenir : $$f(t)=\frac{\sin(t)}{\sqrt{t}}+O\left(\frac{1}{t^{3/2}}\right)$$ En intégrant par parties, on montre que l'intégrale de $\displaystyle\frac{\sin(t)}{\sqrt{t}}$ converge, ce qui permet de conclure.
\end{proof}
\begin{ex}{Une étude d'intégrale semi-convergente}
	On considère l'intégrale $I=\displaystyle\int_{0}^{+\infty}\frac{\sin^3(x)}{x^2}\ dx$ et on définit la fonction $$g:x>0\mapsto\int_{x}^{3x}\frac{\sin(t)}{t^2}\ dt$$
	\begin{enumerate}
		\item Justifier que l'intégrale $I$ est bien convergente. 
		\item Montrer que $g$ est de classe $\mathcal{C}^2$ sur $\R_+\st$ et calculer sa dérivée.
		\item Déterminer la limite de $g$ en $+\infty$ et en $0^+$.
		\item En déduire la valeur de $I$.
	\end{enumerate}
\end{ex}
\begin{proof}
	\begin{enumerate}
		\item Continue sur $\R_+$ après prolongement par continuité en $0$. En $+\infty$, domination par $1/x^2$.
		\item On écrit $g$ à l'aide d'une primitive. En dérivant, on obtient : $$g'(x)=\frac{1}{3x^2}(\sin(3x)-3\sin(x))$$
		\item En $0^+$ : faire un DL de l'intégrande de $g$ et utiliser l'intégration des relations de comparaison pour obtenir $$g(x)=\ln(3)+O\left(x^2\right)$$ donc une limite qui vaut $\ln(3)$ Majorer brutalement en $+\infty$ pour montrer que $g$ tend vers $0$ en $+\infty$.
		\item Utiliser la formule : $$\sin^3(x)=-\frac{1}{4}(\sin(3x)-3\sin(x))$$ On obtient alors en remplaçant dans l'expression de $I$ : $$I=-\frac{3}{4}\int_{0}^{+\infty} g'$$ On obtient donc $$I=\frac{3}{4}\ln(3)$$
	\end{enumerate}
\end{proof}
\section{TD Châteaux}
\begin{ex}{Inégalité de Hardy}
	Soit $f:\R_+\R$ de carré intégrable. On pose, pour tout $x>0$ : $$F(x)=\frac{1}{x}\int_{0}^{x}f(t)\ dt$$ Montrer que $F$ est de carré intégrable et que $\displaystyle\int_{0}^{+\infty}F^2\leqslant 4\int_{0}^{+\infty}f^2$. Indication : on pourra réaliser une intégration par parties.
\end{ex}
\begin{proof}
	Faire une intégration par parties en primitivant le $1/x^2$ puis utiliser l'inégalité de Cauchy-Schwarz, diviser, passer au carré et passer à la limite.
\end{proof}
\begin{ex}{Inégalité de type $L^2$}
	Soit $f:\R\to\R$ de classe $\mathcal{C}^2$ telle que $f^2$ et $(f'')^2$ soient intégrables. 
	\begin{enumerate}
		\item Justifier que $ff''$ est intégrable.
		\item On suppose $(f')^2$ non intégrable sur $\R_+$. Montrer que $ff'$ a pour limite $+\infty$ en $+\infty$ et en déduire une contradiction.
		\item Conclure que $(f')^2$ est intégrable et que $$\left(\int_{-\infty}^{+\infty}(f')^2\right)^2\leqslant\int_{-\infty}^{+\infty}f^2\int_{-\infty}^{+\infty}(f'')^2$$
	\end{enumerate}
\end{ex}
\begin{proof}
	Exercice classique.
	\begin{enumerate}
		\item Classique : utiliser $\lvert ab\rvert\leqslant\displaystyle\frac{1}{2}(a^2+b^2)$.
		\item Intégrer $(ff')'$ entre $0$ et $x$ puis passer à la limite. On en déduit en intégrant désormais $ff'$ que $f^2$ tend vers $+\infty$ en $+\infty$, ce qui contredit l'intégrabilité de $f^2$.
		\item La deuxième question fournit l'intégrabilité sur $[0,+\infty[$. Le raisonnement est symétrique sur $]-\infty,0]$. En intégrant $(ff')'$, on prouve que $ff'$ admet une limite finie, qui ne peut qu'être nulle sans quoi elle contredirait l'intégrabilité de $f^2$. En intégrant $(ff')'$ sur $\R_+$ puis sur $\R_-$ de même, en sommant et par Chasles, on trouve : $$\int_\R(f')^2+\int_\R ff''=0$$ puis on conclut par l'inégalité de Cauchy-Schwarz.
	\end{enumerate}
\end{proof}
\begin{ex}{Inégalité de Wirtinger}
	Soit $f:[0,\pi]\to\R$ de classe $\mathcal{C}^1$ telle que $f(0)=f(\pi)=0$.
	\begin{enumerate}
		\item Justifier l'existence des intégrales ci-dessous et l'égalité : $$\int_{0}^{\pi}f(t)f'(t)\cot(t)\ dt=\frac{1}{2}\int_{0}^{\pi}f(t)^2\left(1+\cot(t)^2\right)\ dt$$
		\item En déduire que $\displaystyle\int_{0}^{\pi}f(t)^2\ dt\leqslant\int_{0}^{\pi}f'(t)^2\ dt$. Étudier le cas d'égalité.
		\item Étendre cette inégalité  à une fonction de $[a,b]$ dans $\R$ telle que $f(a)=f(b)=0$.
	\end{enumerate}
\end{ex}
\begin{proof}
	\begin{enumerate}
		\item Faire une intégration par parties sous réserve d'existence pour l'égalité. Pour l'existence, travailler avec des équivalents.
		\item Partir de $(f(x)\cot(x)-f'(x))^2\geqslant0$, développer et intégrer puis conclure avec la première question. Cas d'égalité lorsque $f'=f\cot$, ce qui donne $f$ proportionnelle à un sinus.
		\item Se ramener à $[0,\pi]$ par translation. On obtient finalement : $$\int_{a}^{b}f^2\leqslant\left(\frac{b-a}{\pi}\right)^2\int_{a}^{b}(f')^2$$
	\end{enumerate}
\end{proof}
\begin{ex}{Lemme de Van Der Corput}
	\begin{enumerate}
		\item Montrer l'existence d'une constante $C>0$ telle que pour toute fonction $\varphi$ de classe $\mathcal{C}^2$ de $[a,b]$ dans $\R$ telle que $\varphi'$ soit monotone et vérifie $\lvert\varphi'\rvert\geqslant1$, on ait : $$\forall\lambda>0,\ \left\lvert\int_{a}^{b}e^{i\lambda\varphi(x)}\ dx\right\rvert\leqslant\frac{C}{\lambda}$$
		\item Soit $k$ un entier supérieur à $2$. Montrer l'existence d'une constante $C_k>0$ telle que pour tout fonction de classe $\mathcal{C}^{k+1}$ de $[a,b]$ dans $\R$ telle que $\lvert\varphi^{(k)}\rvert\geqslant 1$, on ait : $$\forall\lambda>0,\ \left\lvert\int_{a}^{b}e^{i\lambda\varphi(x)}\ dx\right\rvert\leqslant\frac{C_k}{\lambda^{1/k}}$$
	\end{enumerate}
\end{ex}
\begin{proof}
	Écrire $$e^{i\lambda\varphi(x)}=\frac{i\lambda\varphi'(x)}{i\lambda\varphi'(x)}e^{i\lambda\varphi(x)}$$ puis peut-être une IPP ?
\end{proof}
\begin{ex}{Découpage et utilisation des $\varepsilon$}
	Soit $f:\R\to\C$ continue et intégrable sur $\R$.
	\begin{enumerate}
		\item Montrer que $\displaystyle\int_{-\infty}^{+\infty}\left[f(x+t)-f(t)\right]\ dt\xrightarrow[x\to+\infty]{}0$.
		\item Montrer que $\displaystyle\int_{-\infty}^{+\infty}\left\lvert f(x+t)-f(t)\right\rvert\ dt\xrightarrow[x\to+\infty]{}2\int_{-\infty}^{+\infty}\lvert f(t)\rvert\ dt$.
	\end{enumerate}
\end{ex}
\begin{proof}
	Considérer $$S_{X,Y}(x)=\int_{-Y}^{X}(f(x+t)-f(t))\ dt$$ pour $X,Y>0$, séparer en deux intégrales, faire un petit changement de variables et utiliser la relation de Chasles. Finir avec des $\varepsilon$.
\end{proof}
\begin{ex}{Inégalités de Hölder et Minkowski}
	Soient $f$ et $g$ deux fonctions continues sur un segment $I=[a,b]$ à valeurs dans $\K$ et $p$ et $q$ deux exposants conjugués.
	\begin{enumerate}
		\item Démontrer que $\displaystyle\left\lvert\int_{a}^{b}f(t)g(t)\ dt\right\rvert\leqslant\left(\int_{a}^{b}\lvert f(t)\rvert^p\ dt\right)^{1/p}\left(\int_{a}^{b}\lvert g(t)\rvert^q\ dt\right)^{1/q}$.
		\item Démontrer que $\displaystyle\left(\int_{a}^{b}\lvert f(t)+g(t)\rvert^p\ dt\right)^{1/p}\leqslant\left(\int_{a}^{b}\lvert f(t)\rvert^p\ dt\right)^{1/p}+\left(\int_{a}^{b}\lvert g(t)\rvert^p\ dt\right)^{1/p}$.
		\item En déduire que l'application $\lVert\cdot\rVert_p:f\mapsto\displaystyle\left(\int_{a}^{b}\lvert f(t)\rvert^p\ dt\right)^{1/p}$ est une norme sur $\mathcal{C}^0(I,\K)$. Généralise à l'espace $L^p(I,\K)$ des fonctions $f$ continues sur $I$ telles que $\lvert f\rvert^p$ soit intégrable.
		\item Montrer que si $p_1\ne p_2$, les normes $\lVert\cdot\rVert_{p_1}$ et$ \lVert\cdot\rVert_{p_1}$ ne sont pas équivalentes. Trouver une inégalité entre les deux.
		\item Soit $p>1$ un réel et $f:\R_+\to\R_+$ continue telle que $f^p$ soit intégrable. On pose $F(x)=\displaystyle\int_{0}^{x}f(t)\ dt$. Montrer que $x\mapsto\displaystyle\left(\frac{F(x)}{x}\right)^p$ est intégrable et que $$\int_{0}^{+\infty}\left(\frac{F(x)}{x}\right)\ dx\leqslant C_p\int_{0}^{+\infty}\ dx$$ où $C_p$ est une constante à déterminer. C'est l'inégalité de Hardy.
	\end{enumerate}
\end{ex}
\begin{proof}
	Faire comme dans le cas de l'exercice sur les inégalités de Hölder et Minkowski pour les vecteurs de $\K^n$. Pour la dernière, effectuer une intégration par parties puis appliquer l'inégalité de Hölder. On trouve une constante : $$C_p=\left(\frac{p}{p-1}\right)^p$$
\end{proof}

\end{document}